% !TEX root = ../../../tree_notes.tex
% =====================================================================
%  Chapter 12 review -- the 15-section revision contract (chapter_end.txt).
%  Revision-only: tables, one-liners, pointers back via \thmref/\probref/
%  \secref/\coderef/\defref/\exref/\figref/\lemref.  No derivations.
%  Summit: CF 1092F Tree with Maximum Cost -- the weighted distance-sum
%  reroot, f[v] = f[u] + SUM - 2*sz[v], one size/mass DFS + one forward
%  replay.  Sample verified live (Codeforces statement / Luogu mirror):
%  n=8, a=9 4 1 7 10 1 6 5, edges 1-2 2-3 1-4 1-5 5-6 5-7 5-8 -> 121
%  (root 3; profile 2*9+1*4+0*1+3*7+3*10+4*1+4*6+4*5 = 121); n=1 -> 0.
%  Hand-checked here: total=43, f[1]=47, f[3]=f[2]+43-2*sz[3], f[2]=80,
%  f[3]=121 = the maximum.  Exports prob:c09-summit, lst:c09-summit,
%  lst:c09-kit (the last two of this chapter's forward contracts).
% =====================================================================

\chaptersection{Chapter 12 Summary}

One sentence: \emph{a tree's metrics are all read off one unrooted $\Theta(n)$ pass --
$\depth/\height/\size$ from a parent-guarded DFS, the diameter from two BFS sweeps (or a
$\max\max$ of the two longest downward arms at a peak), the centroid from the same size
table via $\mathit{maxpart}$, and every eccentricity from the two diameter ends -- so no
metric query is ever worth more than a couple of linear scans}. Chapter~1 grew height and
size \emph{recursively on binary trees}; this chapter re-frames the identical induction on
\emph{general, unrooted} trees with an explicit parent parameter, and that one change is
what makes the diameter, the centroid and the centre definable at all. Every entry below
is a pointer, not a re-derivation.

\begin{center}
\renewcommand{\arraystretch}{1.25}
\begin{tabular}{@{}p{0.24\linewidth} p{0.48\linewidth} p{0.20\linewidth}@{}}
\toprule
\textbf{Result} & \textbf{Statement (with conditions)} & \textbf{Where} \\
\midrule
Measures & On a general tree rooted at $r$: $\depth$ is pushed \emph{down} (parent $+1$),
  $\height$ and $\size$ folded \emph{up} ($\max_c \height(c){+}1$, $\sum_c \size(c){+}1$) --
  Ch.~1's recursions, now over a child list with a parent guard. & \defref{c09-measures} \\
One pass & A single DFS computes all three arrays in $\Theta(n)$; $\size(r)=n$ is the
  built-in sanity check. & \thmref{c09-onepass} \\
Iterative & Reversed-preorder replay replaces the recursion so a $5\cdot10^{5}$-node chain
  cannot overflow the stack. & \coderef{c09-iter} \\
Diameter & $D$ = \#edges of the longest simple path; two BFS sweeps (farthest-from-anywhere,
  then farthest-from-\emph{that}) reach it, because a farthest node is always peripheral. &
  \defref{c09-diameter}, \thmref{c09-twosweep}, \lemref{c09-sweepend} \\
DP diameter & Equivalently $D = \max_t(\mathit{arm}_1(t)+\mathit{arm}_2(t))$, the two longest
  downward arms joined at the path's peak -- one DFS, survives weights verbatim. & \thmref{c09-dpjoin} \\
Centroid & $c$ with every shard of $T-c$ of size $\le \lfloor n/2\rfloor$; $\mathit{maxpart}$
  re-reads the size table (the upward piece $n-\size(u)$ is the one people forget). & \defref{c09-centroid} \\
Existence & Every tree has a centroid (walk-to-heavy); there are one or two, and when two they
  are adjacent with an exact $n/2{+}n/2$ split. & \thmref{c09-centexist} \\
Find & One DFS decides every node: $c$ is a centroid iff $\mathit{maxpart}(c)\le\lfloor n/2\rfloor$. &
  \thmref{c09-centfind} \\
Eccentricity & $\mathit{ecc}(v)=\max_u d(v,u)$, radius $R=\min_v\mathit{ecc}$, centre the argmin set. &
  \defref{c09-ecc} \\
Two ends suffice & $\mathit{ecc}(v)=\max(d(v,A),d(v,B))$ for \emph{any} diameter ends $A,B$ -- so the
  whole table costs the sweeps you already ran. & \thmref{c09-eccend} \\
Centre rule & The centre is the middle node (even $D$) or middle edge's two nodes (odd $D$) of
  every diameter; $R=\lceil D/2\rceil$; centroid $\ne$ centre. & \thmref{c09-centerrule} \\
\bottomrule
\end{tabular}
\end{center}

\medskip
{\small\textbf{Algorithms from this chapter -- scan the tree a constant number of times, never
per pair.}}
\begin{center}
\small
\renewcommand{\arraystretch}{1.2}
\setlength{\tabcolsep}{3pt}
\begin{tabular}{@{}p{0.19\linewidth} l l p{0.29\linewidth} p{0.155\linewidth}@{}}
\toprule
\textbf{Algorithm} & \textbf{Time} & \textbf{Memory} & \textbf{When to use (deciding phrase)} & \textbf{Code} \\
\midrule
Measures DFS (depth, height, size)
  & $\Theta(n)$ & $\Theta(n)$
  & ``depth / height / subtree size of a general tree''
  & \coderef{c09-dfs3} \\
Iterative two-phase DFS
  & $\Theta(n)$ & $\Theta(n)$
  & ``same, but the chain may hit $5\cdot10^{5}$ deep''
  & \coderef{c09-iter} \\
Two-sweep diameter
  & $\Theta(n)$ & $\Theta(n)$
  & ``the one longest-path length, unweighted''
  & \coderef{c09-sweep} \\
DP diameter (join arms)
  & $\Theta(n)$ & $\Theta(n)$
  & ``diameter that must survive weights / a predicate''
  & \coderef{c09-dp} \\
Centroid scan ($\mathit{maxpart}$)
  & $\Theta(n)$ & $\Theta(n)$
  & ``a safe node to divide-and-conquer on''
  & \coderef{c09-cent} \\
Ecc / radius / centre readout
  & $\Theta(n)$ & $\Theta(n)$
  & ``best worst-case root, or the whole distance profile''
  & \coderef{c09-eccsweep} \\
Distance-sum reroot (Summit)
  & $\Theta(n)$ & $\Theta(n)$
  & ``$\sum_u \mathit{dist}(u,v)\cdot a_v$ for \emph{every} root $v$''
  & \coderef{c09-summit} \\
\bottomrule
\end{tabular}
\end{center}

\chaptersection{Master Decision Flowchart}

Post-study ladder: walk top to bottom, stop at the first branch that matches. It decides
\emph{which single tree metric the question is really asking for} -- and, crucially, hands
off the two questions this chapter does \emph{not} answer: per-pair \emph{query} distance
(Chapter~16's LCA) and the recursion \emph{over} centroids (Chapter~24).

\begin{lstlisting}[language={}, basicstyle=\ttfamily\small, numbers=none, frame=single]
  You are handed an unrooted tree and asked for a number (or a node). Which?
  |
  |-- how deep / how tall / how many nodes below x?
  |     +-- one parent-guarded DFS: depth down, height+size up       [9.1, LC 742]
  |
  |-- the LONGEST path in the tree (its length, maybe its ends)?
  |     +-- two BFS sweeps (unweighted, simplest)                   [9.2, LC 1245]
  |        ...or the DP join of the two longest arms (weights, predicates) [9.2]
  |
  |-- a node whose removal leaves no piece bigger than n/2 (safe to recurse on)?
  |     +-- maxpart = max(n - sz[u], largest child piece) <= n/2    [9.3, CSES 2079]
  |
  |-- the best single root under a WORST-CASE distance (min the farthest)?
  |     +-- ecc[v] = max(d(v,A), d(v,B)) from the two diameter ends  [9.4]
  |        radius = min ecc; centre = the diameter's middle          [9.4]
  |
  |-- the best single root under a SUM of distances (weighted, all roots)?
  |     +-- REROOT SHIFT f[v] = f[u] + SUM - 2*sz[v]                 [9.4, Summit 1092F]
  |        (unweighted count form -> CF 219D's +/-1 reversal cost)   [9.4]
  |
  |-- d(u, v) for MANY arbitrary pairs, online?
  |     +-- NOT this chapter -- LCA arithmetic, Chapter 16 (c11)     [hand off]
  |
  |-- count / combine over ALL paths by recursing the n/2 split?
  |     +-- NOT this chapter -- centroid DECOMPOSITION, Chapter 24 (c18) [hand off]
  |
  +-- a BINARY tree, height/size only, no unrooted metric?
        +-- Chapter 1's recursion already solved it; this chapter is
            the GENERAL-TREE reframing, not a new binary recipe       [cite Ch.1]
\end{lstlisting}

\chaptersection{Common Mistakes Checklist}

\begin{itemize}[nosep, leftmargin=1.4em]
  \item \textbf{Forgetting the upward piece at the centroid.} $\mathit{maxpart}(u)$ is the
        $\max$ of $u$'s \emph{child} subtrees \textbf{and} the remainder $n-\size(u)$ going
        through the parent; drop that last term and a star's centre still passes while a
        chain's does not (\defref{c09-centroid}, \coderef{c09-cent}). This is Trap~1 of
        \secref{c09-centroid} and the single most common submission on CSES~2079.
  \item \textbf{Reporting a single sweep's distance as the diameter.} \coderef{c09-sweep}
        must be \emph{called twice}; one BFS from an arbitrary node returns that node's
        eccentricity, not $D$ -- on the running tree $T$, one sweep from node~4 reports
        $6$, not $D=7$ (\thmref{c09-twosweep}, \probref{c09-diam}).
  \item \textbf{Reading $\le$ vs $<$ at $n/2$.} The centroid bound is $\lfloor n/2\rfloor$,
        attained with \emph{equality} in the two-centroid case; strict ``$< n/2$'' silently
        rejects a valid answer. $T$'s centroids are the adjacent pair $\{1,3\}$, each with a
        piece of \emph{exactly} $n/2$ (\thmref{c09-centexist}, \exref{c09-cent-run}).
  \item \textbf{Calling the centroid the ``centre'' (or vice versa).} The centroid minimises
        the largest \emph{count}; the centre minimises the maximum \emph{distance}. They agree
        on $T$ by luck and disagree on the broom (\secref{c09-radius}).
  \item \textbf{Taking $\mathit{ecc}$ from one diameter end only.} \thmref{c09-eccend} needs
        the $\max$ over \emph{both} ends $A$ and $B$; using one array is the per-source
        distance, not the eccentricity, and understates every node on $A$'s side.
  \item \textbf{Trusting the recursion on a legal-but-deep chain.} \thmref{c09-onepass} costs
        $\Theta(h)$ stack; CSES and most CF constraints allow $h=n$, so ship
        \coderef{c09-iter} whenever $n$ is large -- a stack overflow is a runtime error no
        hand-made sample reveals (\secref{c09-measures}).
  \item \textbf{Missing the sign of the reroot shift.} Crossing $u\to v$: $\size(v)$ nodes get
        \emph{nearer}, $n-\size(v)$ get \emph{farther}, so $f[v]=f[u]+n-2\size(v)$ (weighted:
        $+S-2W_v$). Flip the sign and CF~219D reports the reversal count of the wrong capital
        (\probref{c09-capital}, \probref{c09-summit}).
  \item \textbf{Re-rooting from a size built on the wrong root.} The subtree masses $\size(\cdot)$
        in the reroot replay are those of \emph{one} fixed initial root; recompute them per
        candidate root and you are back to $\Theta(n^2)$ (\coderef{c09-summit}).
\end{itemize}

\chaptersection{WA Clinic}

Verdicts are symptoms; the causes are the short list this chapter drilled. Diagnose
top-down -- every leaf names the training problem.

\begin{figure}[H]
\centering
\begin{tikzpicture}[diagram,
    v/.style={boxnode, fill=primaryblue!12, draw=primaryblue, text width=22mm, align=center},
    b/.style={boxnode, fill=difforange!12, draw=difforange, text width=12mm},
    l/.style={boxnode, fill=successgreen!10, draw=successgreen, text width=58mm,
              font=\scriptsize, align=left}]
  \node[v] (verdict) {\textbf{Verdict on the judge}};
  \node[b, below left=8mm and 6mm of verdict] (wa) {\textbf{WA}};
  \node[b, below right=1mm of verdict] (tle) {\textbf{TLE}};
  \node[b, below left=8mm and 1mm of verdict, xshift=30mm] (re) {\textbf{RE}};
  \node[b, below right=8mm and 6mm of verdict] (mle) {\textbf{MLE}};
  \node[l, below right=1mm and 2mm of wa.south] (wa1)
    {one sweep called as if two, so you submit an eccentricity not $D$ (\probref{c09-diam});
     dropping the $n-\size(u)$ upward piece (\probref{c09-cent}); reading $\mathit{ecc}$ from a
     single diameter end (\thmref{c09-eccend}); a flipped reroot sign (\probref{c09-capital})};
  \node[l, below right=1mm and 2mm of tle.south] (tle1)
    {all-pairs $\Theta(n^2)$ when two sweeps give $D$ in $\Theta(n)$ (\thmref{c09-twosweep});
     re-rooting from scratch per candidate instead of one replay (\probref{c09-summit})};
  \node[l, below right=1mm and 2mm of re.south] (re1)
    {recursive DFS on a legal $5\cdot10^{5}$-node chain -- the stack, not the answer, dies;
     ship \coderef{c09-iter} (\secref{c09-measures})};
  \node[l, below right=1mm and 2mm of mle.south] (mle1)
    {$\Theta(n^2)$ distance matrix materialised for eccentricities you can derive from two
     arrays (\thmref{c09-eccend}); storing a $\mathit{dist}[]$ per sweep when one is reused};
  \draw[flow] (verdict) -- (wa); \draw[flow] (verdict) -- (tle);
  \draw[flow] (verdict) -- (re); \draw[flow] (verdict) -- (mle);
  \draw[flow] (wa) -- (wa1); \draw[flow] (tle) -- (tle1);
  \draw[flow] (re) -- (re1); \draw[flow] (mle) -- (mle1);
\end{tikzpicture}
\caption*{Chapter~12 verdict tree -- only the failures this chapter's problems can produce.}
\end{figure}

\medskip
\noindent Backup table (same content, terse):
\begin{center}
\small
\renewcommand{\arraystretch}{1.2}
\begin{tabular}{@{}llp{0.39\linewidth}l@{}}
\toprule
\textbf{Verdict} & \textbf{Usual cause} & \textbf{One-line fix} & \textbf{Trained in} \\
\midrule
WA  & single sweep called twice once   & \code{D = sweep(sweep(s).far).dist}       & \probref{c09-diam} \\
WA  & no upward piece in maxpart       & \code{max(p, n - sz[u])} before the test   & \probref{c09-cent} \\
WA  & ecc from one end                 & \code{max(dA[v], dB[v])}                   & \thmref{c09-eccend} \\
TLE & all-pairs for $D$                & two BFS sweeps                             & \coderef{c09-sweep} \\
RE  & recursion on a chain             & reversed-preorder replay                   & \coderef{c09-iter} \\
\bottomrule
\end{tabular}
\end{center}

\chaptersection{Quick Reference Card}

\begin{quickref}[Chapter 12 at a glance]
\textbf{One unrooted pass is the budget.} $\depth/\height/\size$ come free from a single
parent-guarded DFS (\thmref{c09-onepass}, \coderef{c09-dfs3}); on a deep chain, swap in the
iterative form \coderef{c09-iter} verbatim -- same three arrays, no stack. Every other metric
here is \emph{that pass, read differently}.

\medskip
\textbf{Two diameter methods, one number.} The sweeps (\coderef{c09-sweep}, unweighted,
shortest code) and the peak-join DP (\coderef{c09-dp}, survives weights and predicates) both
return $D$; the running tree $T$ gives $D=7$ along $8\!-\!5\!-\!2\!-\!1\!-\!3\!-\!7\!-\!9\!-\!10$
(\thmref{c09-twosweep}, \thmref{c09-dpjoin}, \exref{c09-diam-run}).

\medskip
\textbf{Counting vs distance middles.} The centroid balances \emph{head-counts}
($\mathit{maxpart}\le\lfloor n/2\rfloor$, \thmref{c09-centfind}); the centre balances
\emph{mileposts} ($\max(d(\cdot,A),d(\cdot,B))$, \thmref{c09-eccend}). $T$'s centroid pair and
centre pair both happen to be $\{1,3\}$; the broom breaks the coincidence (\secref{c09-radius}).

\medskip
\textbf{Sums reroot; maxima do not.} ``Farthest'' obeys the two-endpoint shortcut
(\thmref{c09-eccend}); ``sum of all distances'' instead obeys the shift
$f[v]=f[u]+S-2W_v$ -- one DFS, one replay, every root (\coderef{c09-eccsweep}), which is the
whole of \probref{c09-capital} and of the Summit \probref{c09-summit}.

\medskip
\textbf{Know your boundaries.} Per-pair query distance is LCA (Chapter~16, \chref{c11});
recursing the $n/2$ split is centroid decomposition (Chapter~24, \chref{c18}); general
rerooting theory is Chapter~15 (\chref{c20}); Euler-tour subtree flattening is Chapter~5
(\chref{c24}). This chapter owns the \emph{one-shot metrics}.
\end{quickref}

\begin{lstlisting}[caption={Chapter 12 code kit: the measures DFS, its iterative rewrite, the two-sweep and DP diameters, the centroid scan, and the eccentricity/reroot readout}, label={lst:c09-kit}]
#include <bits/stdc++.h>
using namespace std;
int n;  vector<vector<int>> adj;

// (1) Measures: depth down, height + size up.  Root anywhere.         [9.1]
vector<int> dep, hgt, sz;
void dfs3(int u, int p) {                       // lst:c09-dfs3
    sz[u] = 1;  hgt[u] = 0;
    for (int v : adj[u]) if (v != p) {
        dep[v] = dep[u] + 1;                    // push down
        dfs3(v, u);
        sz[u]  += sz[v];                        // fold up: size  = sum
        hgt[u]  = max(hgt[u], hgt[v] + 1);      //         height = max+1
    }
}

// (2) Same, no recursion: preorder then reversed replay (chain-safe). [9.1]
void dfs3_iter(int r) {                          // lst:c09-iter
    vector<int> par(n, -1), order;  order.reserve(n);
    vector<int> st{r};
    while (!st.empty()) {                        // phase 1: parent before child
        int u = st.back();  st.pop_back();  order.push_back(u);
        for (int v : adj[u]) if (v != par[u]) { par[v]=u; dep[v]=dep[u]+1; st.push_back(v); }
    }
    for (int i = n-1; i >= 0; --i) {             // phase 2: children first
        int u = order[i];  sz[u]=1;  hgt[u]=0;
        for (int v : adj[u]) if (v != par[u]) { sz[u]+=sz[v]; hgt[u]=max(hgt[u],hgt[v]+1); }
    }
}

// (3) Two sweeps: farthest-from-anywhere, then farthest-from-that.     [9.2]
//     sweep(src) -> {farthest node, its distance}; D = second distance. lst:c09-sweep
pair<int,int> sweep(int src) {
    vector<int> dist(n, -1);  queue<int> q;  q.push(src);  dist[src]=0;
    int far = src;
    while (!q.empty()) { int u=q.front(); q.pop();
        if (dist[u] > dist[far]) far = u;
        for (int v : adj[u]) if (dist[v]==-1) { dist[v]=dist[u]+1; q.push(v); } }
    return {far, dist[far]};
}
int diameter(int s) { auto a=sweep(s); return sweep(a.first).second; }

// (4) DP diameter: join the two longest downward arms at each peak.    [9.2]  lst:c09-dp
int Dp;  void dfs_diam(int u, int p) {
    int a1=0, a2=0;  hgt[u]=0;
    for (int v : adj[u]) if (v!=p) { dfs_diam(v,u); int arm=hgt[v]+1;
        if (arm>=a1){a2=a1;a1=arm;} else if (arm>a2) a2=arm;  hgt[u]=max(hgt[u],arm); }
    Dp = max(Dp, a1 + a2);
}

// (5) Centroid: maxpart with the upward piece, integer-safe test.      [9.3]  lst:c09-cent
int maxpart(int u, int p) {
    sz[u]=1;  int best=0;
    for (int v : adj[u]) if (v!=p) { best=max(best, maxpart(v,u)); sz[u]+=sz[v]; }
    best = max(best, n - sz[u]);                 // THE UPWARD PIECE  <-- do not drop
    return best;                                 // centroid  <=>  2*best <= n
}

// (6) Ecc table + reroot shift: two ends give maxima; one replay gives sums. [9.4] lst:c09-eccsweep
//     ecc[v] = max(dA[v], dB[v])  where A,B = the two diameter ends;
//     sum-costs: f[0]=sum dep[v]*a[v], then f[v]=f[u]+total-2*sz[v] over each tree edge u->v.
\end{lstlisting}

\chaptersection{Write-From-Memory Drills}

Five signature skeletons, each lifted from a listing already in the chapter. Cover the
marked line, reproduce it from memory, then check the answer one line below.

\begin{lstlisting}[caption={Drill 1 -- the measures DFS: down-push, up-fold (from \coderef{c09-dfs3})}]
void dfs3(int u, int p) {
    sz[u] = 1;  hgt[u] = 0;
    for (int v : adj[u]) if (v != p) {        // parent guard = the rooting
        dep[v] = dep[u] + 1;                   /* FILL: dfs3(v, u); */
        sz[u]  += sz[v];                        // up-fold:  size  = sum of children
        hgt[u]  = max(hgt[u], hgt[v] + 1);     //           height = max child + 1
    }
}
\end{lstlisting}
\noindent{\scriptsize Key: depth is pushed \emph{before} the recursive call, height and size
are folded \emph{after} it -- ``depth down, height/size up'' (Ch.~1's rule, now over a child
list). The guard ``\code{v != p}'' is what lets one DFS serve every root choice
(\defref{c09-measures}, \thmref{c09-onepass}).}

\begin{lstlisting}[caption={Drill 2 -- the chain-safe two-phase DFS (from \coderef{c09-iter})}]
while (!st.empty()) { int u = st.back(); st.pop_back(); order.push_back(u);
    for (int v : adj[u]) if (v != par[u]) { par[v]=u; dep[v]=dep[u]+1; /* FILL: st.push_back(v); */ } }
for (int i = n-1; i >= 0; --i) {              // reversed preorder == children first
    int u = order[i];  sz[u]=1;  hgt[u]=0;
    for (int v : adj[u]) if (v != par[u]) { sz[u]+=sz[v]; hgt[u]=max(hgt[u],hgt[v]+1); }
}
\end{lstlisting}
\noindent{\scriptsize Key: a \emph{tree} preorder written parent-before-child, read backwards,
is a valid bottom-up order -- so the recursion becomes two loops and one stack vector. This is
the rewrite Chapter~16's kit cites for $5\cdot10^{5}$-node chains (\thmref{c09-onepass}).}

\begin{lstlisting}[caption={Drill 3 -- the two sweeps (from \coderef{c09-sweep})}]
int diameter(int s) {
    auto A = sweep(s);             // farthest from anywhere  -> an END of a diameter
    auto B = sweep(A.first);       // farthest from that END  -> the OTHER end
    /* FILL: return B.second; */   // = D
}
\end{lstlisting}
\noindent{\scriptsize Key: one BFS/DFS distance pass called twice. A single call returns the
\emph{eccentricity of the start}, not $D$; the second call is the whole trick
(\thmref{c09-twosweep}, \lemref{c09-sweepend}).}

\begin{lstlisting}[caption={Drill 4 -- centroid maxpart (from \coderef{c09-cent})}]
int maxpart(int u, int p) {
    sz[u] = 1;  int best = 0;
    for (int v : adj[u]) if (v != p) { best = max(best, /* FILL: maxpart(v, u) */); sz[u] += sz[v]; }
    best = max(best, n - sz[u]);               // the upward piece -- the line people forget
    return best;                                // centroid  <=>  2 * best <= n
}
\end{lstlisting}
\noindent{\scriptsize Key: the shards created by deleting $u$ are its child subtrees \emph{and}
the $n-\size(u)$ remainder through the parent; the integer-safe test \code{2*best <= n} avoids
the $\lfloor n/2\rfloor$ floor and handles the two-centroid equality case (\thmref{c09-centfind}).}

\begin{lstlisting}[caption={Drill 5 -- the reroot shift for sum-of-distances (from \coderef{c09-summit})}]
void dfs2(int u, int p) {                       // forward replay, after sz[] from dfs1
    for (int v : adj[u]) if (v != p) {
        f[v] = f[u] + total - 2 * sz[v];        /* FILL: which term is the mass left behind? */
        dfs2(v, u);
    }
}
\end{lstlisting}
\noindent{\scriptsize Key: crossing edge $u\to v$, the $\size(v)$ nodes inside $v$'s side move
one step \emph{nearer} ($-2\size(v)$ against the all-farther baseline) and the other
$n-\size(v)$ move one step \emph{farther}; \code{total} is $\sum a$ over \emph{all} nodes.
Weighted form replaces counts by mass $W_v$ vs $S-W_v$ (\thmref{c09-eccend} has no $\sum$
analogue -- this is the sum version, \probref{c09-summit}).}

\chaptersection{Complete Problem Checklist}

Every problem this chapter offers -- 4 fully solved in-section, 5 practice, one Summit. Tick
\emph{only} after an unaided solve; a hinted attempt counts as half. The ten mixed problems
are listed in their own station. Cite by name and number -- URLs appear once, in the solved
and practice entries.

\begin{center}
\footnotesize
\renewcommand{\arraystretch}{1.25}
\begin{tabular}{@{}c p{0.19\linewidth} l p{0.20\linewidth} p{0.235\linewidth} c@{}}
\toprule
\textbf{\#} & \textbf{Problem} & \textbf{Source} & \textbf{Key technique} & \textbf{What to look for} & \textbf{Done?} \\
\midrule
\multicolumn{6}{l}{\textbf{Section 9.1 -- depth, height, size on a general tree}} \\
W12.1.1 & Closest Leaf in a Binary Tree & LC 742 & \lv{3} undirected BFS by layers & rebuild as adjacency, BFS from \code{k}, first leaf wins; ``leaf'' is a binary-tree notion & $\square$ \\
P12.1.2 & Evaluate Boolean Binary Tree & LC 2331 & \lv{2} bottom-up fold & leaves \code{0/1}, internal \code{2}=OR \code{3}=AND -- read the encoding, it is the easy trap & $\square$ \\
P12.1.3 & Even Odd Tree & LC 1609 & \lv{3} level-order check & even levels odd \& increasing, odd levels even \& decreasing -- one BFS, per-level compare & $\square$ \\
\midrule
\multicolumn{6}{l}{\textbf{Section 9.2 -- the diameter, two ways}} \\
W12.2.1 & Tree Diameter & LC 1245 & \lv{3} two sweeps (edge count) & answer is the \#edges; any start, BFS twice, report the second distance & $\square$ \\
P12.2.2 & Fixed-Length Paths II & CSES 2081 & \lv{5} peak-join with a count & count pairs with $a\le$ path $\le b$: the DP geometry, $\max$ replaced by a frequency scan & $\square$ \\
\midrule
\multicolumn{6}{l}{\textbf{Section 9.3 -- the centroid}} \\
W12.3.1 & Finding a Centroid & CSES 2079 & \lv{3} maxpart scan & print any $c$ with $\mathit{maxpart}\le n/2$; the upward piece decides pass/fail & $\square$ \\
P12.3.2 & Xenia and Tree & CF 342E & \lv{6} centroid decomp (gated) & red-node nearest query -- \emph{defer} until Ch.~24; here only spot the centroid need & $\square$ \\
\midrule
\multicolumn{6}{l}{\textbf{Section 9.4 -- eccentricity, radius, centre, reroot}} \\
W12.4.1 & Choosing Capital for Treeland & CF 219D & \lv{4} reroot on reversal cost & $f[v]=f[u]+1-2w_{\text{in}}[v]$; print the min-cost city and that cost & $\square$ \\
P12.4.2 & Shortest Distance from All Buildings & LC 317 & \lv{5} multi-source BFS sums & grid, not tree: BFS from each house, sum over cells reaching all; the $\sum$-family sibling & $\square$ \\
\midrule
\multicolumn{6}{l}{\textbf{Summit (chapter end)}} \\
S12 & Tree with Maximum Cost & CF 1092F & \lv{5} weighted distance-sum reroot & $f[v]=f[u]+S-2W_v$; \code{long long}; the max over all roots & $\square$ \\
\bottomrule
\end{tabular}
\end{center}

\chaptersection{Recommended Study Order}

Five focused days. Day~1 rebuilds the three measures on a general tree and drills the
chain-safe rewrite; day~2 finds the diameter two independent ways; day~3 nails the centroid
and the $n/2$ boundary; day~4 reads eccentricity/radius/centre off the diameter ends and
meets the reroot shift; day~5 is the Summit. Solve in this order -- each day's tool is the
previous day's output.

\begin{center}
\small
\renewcommand{\arraystretch}{1.25}
\begin{tabular}{@{}llp{0.55\linewidth}@{}}
\toprule
\textbf{Day} & \textbf{Focus} & \textbf{Tasks} \\
\midrule
1 & Read + do 9.1 & trace $T$'s $\depth/\height/\size$ from root~1 (\exref{c09-measure-run});
    write \coderef{c09-dfs3} then \coderef{c09-iter} cold; LC 742, then 2331, 1609 \\
2 & Read + do 9.2 & get $D=7$ on $T$ by both methods (\exref{c09-diam-run}); prove
    \lemref{c09-sweepend} to yourself; LC 1245, then CSES 2081 \\
3 & Read + do 9.3 & compute $\mathit{maxpart}$ on all of $T$, see the pair $\{1,3\}$
    (\exref{c09-cent-run}); CSES 2079; only \emph{read} CF 342E's statement \\
4 & Read + do 9.4 & fill $T$'s $\mathit{ecc}$ table from ends $8,10$ (\exref{c09-ecc-run});
    derive the shift; CF 219D, then LC 317 \\
5 & Summit + revise & CF 1092F with the $\pm$ reroot replay (\probref{c09-summit}); walk the
    decision flowchart on five questions \\
\bottomrule
\end{tabular}
\end{center}

\chaptersection{Difficulty Progression}

Tiers are competency gates -- each row names the \emph{skill} the tier unlocks, not just a
rating bucket.

\begin{center}
\small
\renewcommand{\arraystretch}{1.3}
\begin{tabular}{@{}lp{0.36\linewidth}p{0.34\linewidth}@{}}
\toprule
\textbf{Tier} & \textbf{Problems} & \textbf{What passing this tier proves} \\
\midrule
$\lv{2}$ \ $\bigstar\bigstar$
  & LC 2331, LC 733
  & you fold a value up a tree with a child list and read the encoding without re-deriving it \\
$\lv{3}$ \ $\bigstar\bigstar\bigstar$
  & LC 742, 1609, 1245, CSES 2079, 1130, LC 2316, 542, 1162
  & you compute size/height in one pass and turn ``longest path'' or ``a balanced cut'' into two scans \\
$\lv{4}$ \ $\bigstar\bigstar\bigstar\bigstar$
  & CF 219D, 337D, LC 2858, 2421
  & you carry a per-node quantity down and replay it across every root with a signed shift \\
$\lv{5}$ \ $\bigstar\bigstar\bigstar\bigstar\bigstar$
  & CSES 2081, LC 317, 1948, CF 592D, CF 1092F (Summit)
  & you combine the diameter geometry with counts/weights and choose $\max$ vs $\sum$ deliberately \\
$\lv{6}$ \ $\bigstar\bigstar\bigstar\bigstar\bigstar\bigstar$
  & CF 342E
  & you recognise centroid decomposition as the recursion this chapter only previews (Ch.~24) \\
\bottomrule
\end{tabular}
\end{center}

\chaptersection{Reduction Gallery}

Every metric is ``scan the tree a constant number of times, then read a per-node number.''
Each line collapses a family onto a listing already in the chapter -- no new theory.

\begin{itemize}[nosep, leftmargin=1.4em]
  \item ``closest leaf to node $k$'' (LC 742) $\longrightarrow$ unroot the tree, BFS by layers
        from $k$, first leaf is the answer (\probref{c09-deep}).
  \item ``evaluate a boolean tree bottom-up'' (LC 2331) $\longrightarrow$ one fold at the node,
        reading \code{2}/\code{3} as OR/AND (Ch.~1's recursion, general children).
  \item ``longest path in the tree'' (LC 1245) $\longrightarrow$ two BFS sweeps, second distance
        is $D$ (\thmref{c09-twosweep}, \probref{c09-diam}).
  \item ``count pairs at path length in $[a,b]$'' (CSES 2081) $\longrightarrow$ the peak-join
        geometry of \thmref{c09-dpjoin} with $\max$ replaced by a frequency scan (\probref{c09-diam}'s
        harder sibling).
  \item ``a node whose removal balances the tree'' (CSES 2079) $\longrightarrow$ $\mathit{maxpart}$
        from the size table (\thmref{c09-centfind}, \probref{c09-cent}).
  \item ``best worst-case root / the radius'' (LC 310's centre) $\longrightarrow$
        $\mathit{ecc}(v)=\max(d(v,A),d(v,B))$, argmin is the middle of the diameter
        (\thmref{c09-eccend}, \thmref{c09-centerrule}).
  \item ``capital minimising reversed roads'' (CF 219D) $\longrightarrow$ reroot with edge costs
        $\pm1$: $f[v]=f[u]+1-2w_{\text{in}}[v]$ (\probref{c09-capital}).
  \item ``maximum weighted distance sum over every root'' (CF 1092F) $\longrightarrow$ the same
        replay with mass: $f[v]=f[u]+S-2W_v$ (\probref{c09-summit}).
\end{itemize}

\chaptersection{Mixed Problem Set}

These refuse to say ``metric''; decide \emph{which quantity lives at each node and whether the
answer is a $\max$, a $\sum$, or a count} before opening the code. A hint tells you the shape
of the answer, never the tool.

\begin{mixedset}
  \practice{Super M}{Codeforces 592D}
    {https://codeforces.com/problemset/problem/592/D}{CF 2200 · L5}
    {Cows visit $m$ marked nodes of a tree and return; the last cow may stop anywhere. Minimise the
    total traversed edges. Read the answer off the diameter of the \emph{marked} subtree: it is
    $2E - D$ where $E$ is that subtree's edge count and $D$ its diameter -- start at the lower-numbered
    end of the marked-diameter. Exactly \thmref{c09-twosweep} run on a pruned tree (\probref{c09-diam}).}
  \practice{Book of Evil}{Codeforces 337D}
    {https://codeforces.com/problemset/problem/337/D}{CF 1800 · L4}
    {Place one root so that every cursed node is within distance $d$ of it; count the valid nodes. Two
    DFS passes bound each node by its farthest cursed node above and below -- a per-node eccentricity
    restricted to a marked set, i.e.\ \thmref{c09-eccend} with a maximum-distance-to-markes table
    instead of to diameter ends.}
  \practice{Tree Matching}{CSES 1130}
    {https://cses.fi/problemset/task/1130}{CSES · L4}
    {Maximum set of vertex-disjoint edges. The bottom-up fold ``take the edge to a free child or don't''
    is Ch.~13's tree-DP template, but the greedy ``match a leaf to its parent, delete both, repeat''
    needs only the size/height discipline of \secref{c09-measures} to justify its $\Theta(n)$ implementation.}
  \practice{Minimum Edge Reversals So Every Node Is Reachable}{LeetCode 2858}
    {https://leetcode.com/problems/minimum-edge-reversals-so-every-node-is-reachable/}{LC Hard · L4}
    {Directed tree (an orientation of a spanning tree); for each node $i$ count the fewest edges to flip
    so all others are reachable from $i$. This is \probref{c09-capital}'s reversal cost asked for every
    root: one DFS for $\mathit{ans}[0]$, then $ans[v]=ans[u]+1-2\cdot[\text{edge } u\!\to\!v \text{ exists}]$.}
  \practice{Number of Good Paths}{LeetCode 2421}
    {https://leetcode.com/problems/number-of-good-paths/}{LC Hard · L4}
    {A ``good path'' has endpoints of equal value and no intermediate larger. Process values ascending,
    union-find on edges whose endpoints are $\le$ the current value; each component of size $s$
    contributes $s(s+1)/2$. The $\mathit{size}[]$ you maintain is exactly the $\size$ of \secref{c09-measures}
    re-purposed for counting.}
  \practice{Critical Connections in a Network}{LeetCode 1948}
    {https://leetcode.com/problems/critical-connections-in-a-network/}{LC Hard · L5}
    {Remove single connections to split the network (bridges) and find duplicate roads. Tarjan's
    bridge-finding rides on the $\depth$ / low-link distinction of \secref{c09-measures}: a tree edge
    $u\!\to\!v$ is a bridge iff nothing on $v$'s side reaches above $v$. (The general-graph DFS is
    Chapter~29's; the rooted $\depth$ vocabulary is this chapter's.)}
  \practice{Count Unreachable Pairs of Nodes}{LeetCode 2316}
    {https://leetcode.com/problems/count-unreachable-pairs-of-nodes-in-an-undirected-graph/}{LC Medium · L3}
    {Not a tree, but each DFS component yields a $\size$; the unreachable pairs are $\sum$ over components
    of $\size\cdot(\text{already seen})$. One size table (\thmref{c09-onepass}) is the whole arithmetic.}
  \practice{Flood Fill}{LeetCode 733}
    {https://leetcode.com/problems/flood-fill/}{LC Easy · L2}
    {Recolour the 4-connected region of a pixel. The BFS/DFS ``visit every reachable node'' that powers
    \coderef{c09-dfs3} is the entire algorithm on a grid; the running tree's phase-1 preorder is the
    same wavefront.}
  \practice{01 Matrix}{LeetCode 542}
    {https://leetcode.com/problems/01-matrix/}{LC Medium · L3}
    {Distance of each cell to the nearest \code{0}. Multi-source BFS from all zeros at once -- the
    ``pop order is distance order'' fact behind \lemref{c09-sweepend}, with every source at depth~0.}
  \practice{As Far from Land as Possible}{LeetCode 1162}
    {https://leetcode.com/problems/as-far-from-land-as-possible/}{LC Medium · L3}
    {Find the water cell maximising its distance to the nearest land: the multi-source-BFS eccentricity
    of the \emph{grid}, i.e.\ $\max_v \min_{\text{land } \ell} d(v,\ell)$ -- the mirror image of the
    radius of \defref{c09-ecc} with $\max$ and $\min$ swapped.}
\end{mixedset}

\chaptersection{The Summit -- CF 1092F Tree with Maximum Cost}

\begin{notebox}[Why this is the chapter top rope]
Sections 9.1--9.4 all measured a tree with \emph{one} root chosen for you, or with the two
sweeps' distance arrays. The Summit breaks that assumption cleanly: you get to \emph{pick the
root} $v$ and must maximise a weighted distance sum $\sum_u \mathit{dist}(u,v)\cdot a_u$ over
all $n$ choices. Naively that is $n$ DFS passes -- $\Theta(n^2)$ at $n=2\cdot10^{5}$. The fix
is the reroot shift this chapter previewed at \probref{c09-capital}: the answer for a child is
the answer for its parent plus a correction that depends only on the mass on each side of the
edge, so one size/mass DFS plus one forward replay returns \emph{every} root's total in
$\Theta(n)$. This is CF~219D with the $\pm1$ edge cost replaced by a weighted mass -- the
\emph{sum} analogue of \thmref{c09-eccend}, and the exact technique Chapter~15 (\chref{c20})
later builds a whole theory on. It is a summit, not a new country: every tool is already in
\coderef{c09-kit}.
\end{notebox}

\solvedproblem{c09-summit}{Tree with Maximum Cost}{Codeforces 1092F}
  {https://codeforces.com/problemset/problem/1092/F}{CF 1900 · L5}
  {weighted distance-sum rerooting, subtree-mass DFS, one forward replay of the parent answer using total minus twice the child mass, long long to survive large sums}

\begin{problemstatement}[Tree with Maximum Cost]
You are given a tree of $n$ vertices; vertex $v$ carries a value $a_v$. For a chosen root $v$,
the \term{cost} is $\sum_{i=1}^{n} \mathit{dist}(i,v)\cdot a_i$, where $\mathit{dist}$ is the
number of edges on the path. Choose $v$ to \emph{maximise} the cost, and print that maximum.\\
$1 \le n \le 2\cdot10^{5}$, $1 \le a_v \le 2\cdot10^{5}$; edges are vertex pairs forming a tree.
\textbf{Example~1}: $n=8$, values \texttt{9 4 1 7 10 1 6 5}, edges \texttt{1-2, 2-3, 1-4, 1-5,
5-6, 5-7, 5-8}. Choosing $v=3$ costs $2\cdot9+1\cdot4+0\cdot1+3\cdot7+3\cdot10+4\cdot1+4\cdot6+4\cdot5
= 121$, and no root beats it $\Rightarrow$ output \texttt{121}. \textbf{Example~2}: $n=1$, value
\texttt{1337} $\Rightarrow$ the only distance is~$0$, output \texttt{0}.
\end{problemstatement}

\begin{attemptgate}[Attempt it yourself first -- \lv{5}, give yourself 40 minutes]
A fair first attempt is the honest one: DFS from each vertex $v$ to sum $\mathit{dist}(v,\cdot)\cdot
a$, take the max. It is correct and it is $\Theta(n^2)$ -- at $n=2\cdot10^{5}$ that is $4\cdot10^{10}$
and it \tle s. The question is whether the $n$ answers share enough structure to be swept in one pass.
\peek{When the root moves across one edge $u\to v$, how does the cost change -- and which two counts
fully describe that change?}
\nudge{Every node inside $v$'s side of the edge gets one step \emph{nearer}; every node on the other
side gets one step \emph{farther}. So you only need the total value mass on each side -- $f[v]=f[u]
+(\text{mass outside}) - (\text{mass inside})$.}
\fullhint{Root anywhere (say node~0). A DFS computes, for each node, its subtree mass $\mathit{sz}[v]
=\sum a$ in $v$'s subtree, and $f[0]=\sum \depth(u)\cdot a_u$ in the same pass. Let $S=\sum a_u$.
Then a forward replay gives $f[v]=f[u]+S-2\,\mathit{sz}[v]$ for every tree edge $u\to v$ ($u$ the
parent). Print $\max_v f[v]$. Use \code{long long}: $f$ can reach $\sim 2\cdot10^{5}\cdot2\cdot10^{5}$.
Full build below.}
\end{attemptgate}

\solstep{Understanding the Problem}
The objective is a weighted \emph{sum} of distances to a chosen root, and we choose the root to
\emph{maximise} it (note: opposite of a facility-location ``best root'', which minimises -- same
machinery, opposite sign of intent). Two facts decide the algorithm. First, the answer for a
neighbouring root differs from the current root by a quantity computable from two numbers: how
much value sits on each side of the connecting edge. Second, ``how much value is on the child's
side'' is exactly a subtree mass from one rooted DFS. So all $n$ answers are reachable by a walk
that reuses $f[\mathit{parent}]$, and the $\Theta(n^2)$ collapses to $\Theta(n)$.

\solstep{Key Observations}
\begin{observation}
Moving the root across edge $u\to v$ (with $v$ the child under $u$): the $\mathit{sz}[v]$ value
mass in $v$'s subtree moves one step closer ($-\mathit{sz}[v]$ to the cost) and the remaining
$S-\mathit{sz}[v]$ moves one step farther ($+(S-\mathit{sz}[v])$). Net:
$f[v]=f[u]+S-2\,\mathit{sz}[v]$. This is \thmref{c09-onepass}'s size table read as weights.
\end{observation}
\begin{observation}
$f[\mathit{root}]$ and every $\mathit{sz}[\cdot]$ come from \emph{one} DFS rooted arbitrarily --
no per-candidate recomputation. The subtree masses are taken once, w.r.t.\ that fixed root, and
are exactly what the replay needs; re-rooting does not recompute them (\probref{c09-capital} uses
the identical discipline with $\pm1$ instead of mass).
\end{observation}
\begin{observation}
Every partial sum is $\le n\cdot\max a$ and the answer itself can hit $\sim 4\cdot10^{10}$, so
\code{int} overflows -- the \emph{implementation} bug here is a type, not an algorithm. Use
\code{long long} for $a$, $\mathit{sz}$, $S$ and $f$ (\secref{c09-radius}'s reroot preview warned).
\end{observation}

\solstep{Think It Through With Me}
\begin{thinkbox}
\deadend{DFS from every vertex and take the max.}
{Correct but $\Theta(n^2)$; at $n=2\cdot10^{5}$ it is $4\cdot10^{10}$ edge-visits -- \tle. The $n$
answers must be made to \emph{share} work, not be computed independently.}
\deadend{Treat it like the radius and reuse $\max(d(v,A),d(v,B))$.}
{That identity is for a \emph{maximum} over distances (\thmref{c09-eccend}); this objective is a
\emph{weighted sum}, which has no two-endpoint shortcut -- the sums obey a different law (a shift),
not the same $\max$.}
\aha{Compute the answer for \emph{one} root and a subtree-mass table in one DFS, then push answers
to children with $f[v]=f[u]+S-2\mathit{sz}[v]$. Two DFS passes, all roots done, and it is CF~219D's
skeleton with masses replacing the $-1/+1$ reversal counts.}
\end{thinkbox}

\solstep{The Approach Ladder}
\begin{center}\renewcommand{\arraystretch}{1.3}
\begin{tabular}{@{}c p{0.235\linewidth} p{0.30\linewidth} l c@{}}
\toprule
\textbf{\#} & \textbf{Approach} & \textbf{Idea} & \textbf{Time} & \textbf{Verdict} \\
\midrule
1 & \lv{3} DFS from each root & recompute $\sum \mathit{dist}\cdot a$ per vertex & $O(n^2)$ & \tle \\
2 & \lv{4} BFS layers per root & same sum, BFS instead of DFS & $O(n^2)$ & \tle \\
3 & \lv{5} reroot replay & one mass DFS + $f[v]=f[u]+S-2\mathit{sz}[v]$ & $O(n)$ & \ac \\
\bottomrule
\end{tabular}\end{center}
Rungs 1--2 differ only in traversal but share the fatal $n$ factor. Rung 3 is \emph{this} chapter's
answer -- recognise that all roots' totals are one replay apart, exactly as \secref{c09-radius}
previewed and \probref{c09-capital} proved on unit costs.

\solstep{Solution}
Root the tree at node~0. A single DFS computes, for each node $u$, its subtree mass $\mathit{sz}[u]$
(the sum of $a$ over $u$'s rooted subtree) and, carrying $\depth$, the seed cost $f[0]=\sum_u
\depth(u)\,a_u$. Let $S=\sum_u a_u$ be the total mass, computed while reading input. A second
DFS walks downward: for each tree edge $u\to v$ with $u$ the parent, the shift
$f[v]=f[u]+S-2\mathit{sz}[v]$ transfers the known parent answer to the child in $\bigO{1}$ --
$\mathit{sz}[v]$ nodes come nearer, $S-\mathit{sz}[v]$ go farther. The answer is $\max_v f[v]$.
All accumulators are \code{long long}.

\solstep{Proof of Correctness}
\begin{proofbox}[The replay returns the exact cost for every root]
\emph{Seed.} $f[0]$ is by definition $\sum_u \depth_0(u)\,a_u$ where $\depth_0$ is distance from
node~0, computed correctly by the first DFS. \emph{Shift.} Fix a tree edge $u\text{--}v$ with $v$
a child of $u$ under root~0. Removing the edge splits the vertices into $v$'s subtree (mass
$\mathit{sz}[v]$) and the complement (mass $S-\mathit{sz}[v]$). Measuring distance from $v$ rather
than from $u$ decreases the distance to each subtree vertex by exactly one and increases the
distance to each complement vertex by exactly one, so
$f[v]=f[u]-\mathit{sz}[v]+(S-\mathit{sz}[v])=f[u]+S-2\mathit{sz}[v]$. \emph{Coverage.} The second
DFS visits every vertex exactly once with its parent already settled, so every $f[\cdot]$ is
computed from its true parent value; by induction on rooted depth all $f[v]$ are correct. Taking
$\max_v f[v]$ therefore returns the maximum cost, achieved by \thmref{c09-onepass}'s linear bound
on each of the two passes.
\end{proofbox}

\solstep{Implementation}
\begin{lstlisting}[caption={Tree with Maximum Cost -- one subtree-mass DFS + one forward reroot replay f[v]=f[u]+S-2*sz[v]; long long throughout (Codeforces 1092F; O(n))}, label={lst:c09-summit}]
#include <bits/stdc++.h>
using namespace std;
int n;
vector<vector<int>> adj;
vector<long long> a, sz, f, dep;
long long total = 0;
void dfs1(int u, int p) {                 // mass + seed cost f[0] in one pass
    sz[u] = a[u];
    for (int v : adj[u]) if (v != p) {
        dep[v] = dep[u] + 1;
        f[0] += dep[v] * a[v];            // contribution of v measured from root 0
        dfs1(v, u);
        sz[u] += sz[v];
    }
}
void dfs2(int u, int p) {                 // forward replay: parent answer -> child answer
    for (int v : adj[u]) if (v != p) {
        f[v] = f[u] + total - 2 * sz[v];  // sz[v] nearer, (total-sz[v]) farther
        dfs2(v, u);
    }
}
int main() {
    if (scanf("%d", &n) != 1) return 0;
    adj.assign(n, {}); a.assign(n, 0); sz.assign(n, 0); f.assign(n, 0); dep.assign(n, 0);
    for (int i = 0; i < n; i++) { long long x; scanf("%lld", &x); a[i] = x; total += x; }
    for (int i = 0; i < n - 1; i++) { int u, v; scanf("%d %d", &u, &v); --u; --v;
        adj[u].push_back(v); adj[v].push_back(u); }
    dfs1(0, -1);                          // fills sz[], dep[], f[0]
    dfs2(0, -1);                          // fills f[1..n-1]
    long long ans = *max_element(f.begin(), f.end());
    printf("%lld\n", ans);
    return 0;
}
\end{lstlisting}
\begin{notebox}[The one line that is the whole idea]
\code{f[v] = f[u] + total - 2*sz[v];} --- everything else is reading input and taking a max.
Swap \code{sz} for ``number of nodes'' and \code{total} for $n$ and it is the unweighted
distance-sum reroot; set the edge cost to the reversal $\pm1$ and it is \probref{c09-capital}.
\end{notebox}

\solstep{Trace}
\begin{tracebox}[Official Sample 1 -- values 9 4 1 7 10 1 6 5, total = 43, rooted at vertex 1]
Tree under root~1: children $1\!\to\!\{2,4,5\}$, $2\!\to\!3$, $5\!\to\!\{6,7,8\}$.
\begin{itemize}[nosep,leftmargin=1.3em]
  \item \textbf{Masses} $\mathit{sz}$ (bottom-up): $\mathit{sz}[3]=1$, $\mathit{sz}[2]=4+1=5$,
        $\mathit{sz}[4]=7$, $\mathit{sz}[6]=1$, $\mathit{sz}[7]=6$, $\mathit{sz}[8]=5$,
        $\mathit{sz}[5]=10{+}1{+}6{+}5=22$, $\mathit{sz}[1]=9{+}5{+}7{+}22=43$ ($=\mathit{total}$
        \checkmark).
  \item \textbf{Seed} $f[1]=\sum \depth\cdot a$ (depths $0,1,2,1,1,2,2,2$ for $1..8$):
        $0\cdot9+1\cdot4+2\cdot1+1\cdot7+1\cdot10+2\cdot1+2\cdot6+2\cdot5 = 47$.
  \item \textbf{Replay} $f[v]=f[u]+43-2\mathit{sz}[v]$:
        $f[2]=47+43-2{\cdot}5=80$; \ $f[3]=80+43-2{\cdot}1=121$; \ $f[4]=47+43-2{\cdot}7=76$;
        $f[5]=47+43-2{\cdot}22=46$; \ $f[6]=46+43-2{\cdot}1=87$; \ $f[7]=46+43-2{\cdot}6=77$;
        $f[8]=46+43-2{\cdot}5=79$.
\end{itemize}
Values $f=[47,80,\mathbf{121},76,46,87,77,79]$; the maximum is $f[3]=\mathbf{121}$ \checkmark ---
matching the official note (root $3$). \textbf{Sample~2}: $n=1$, no edges, $f[1]=0$, output
\texttt{0} \checkmark.
\end{tracebox}

\solstep{Complexity}
\begin{complexitybox}
Time: two DFS passes, each visiting every edge twice -- $\Theta(n)$ total, versus the
$\Theta(n^2)$ brute force; the shift's $\bigO{1}$ per edge is what buys the linear bound
(\thmref{c09-onepass}). Memory: four $\bigO{n}$ \code{long long} arrays plus the adjacency
$\bigO{n}$. On a legal $n=2\cdot10^{5}$ chain the recursion is $2\cdot10^{5}$ frames deep --
swap \code{dfs1}/\code{dfs2} for the \coderef{c09-iter} two-phase form if your stack is small
(Chapter~16's kit makes exactly this substitution).
\end{complexitybox}

\solstep{Verification and Edge Cases}
\begin{center}\renewcommand{\arraystretch}{1.2}
\begin{tabular}{@{}p{0.30\linewidth} p{0.62\linewidth}@{}}
\toprule
\textbf{Case} & \textbf{What the code does} \\
\midrule
Single node ($n=1$) & no edges, $f[0]=0$, answer \texttt{0} -- Sample~2 \checkmark \\
Star, one heavy leaf & the leaf's $f$ gets $+S-2a_{\text{leaf}}$, so a light root wins -- max picks a leaf side \\
Long chain, all $a=1$ & answer is the sum of distances from an endpoint $=n(n-1)/2$; needs \code{long long} ($n=2\cdot10^{5}$) \\
Root-0 mass check & $\mathit{sz}[\text{root}]=\mathit{total}$ is a free correctness assertion after \code{dfs1} \\
Ties on the max & any argmax is fine -- the problem asks only for the maximum cost \\
\bottomrule
\end{tabular}\end{center}

\solstep{Pitfalls and Variants}
\begin{warningbox}
Ways to break it: (i) \emph{\code{int} overflow} -- $f$ reaches $\sim4\cdot10^{10}$; every
accumulator must be \code{long long}. (ii) \emph{recomputing} $\mathit{sz}$ per candidate root --
the masses are taken once under a fixed root; the replay reuses them, and recomputing is the
$\Theta(n^2)$ \tle. (iii) \emph{sign flip} -- $+S-2\mathit{sz}[v]$, not $+2\mathit{sz}[v]-S$; on
Sample~1 the wrong sign sends $f[3]$ off~121. (iv) \emph{recursion depth} on a chain -- see the
complexity note. \emph{Variants:} set $a\equiv1$ and this is the plain distance-sum reroot;
replace the shift's mass by the reversal $\pm1$ and it is \probref{c09-capital}; the general
``compute $f$ for all roots from an edge-transfer rule'' family is Chapter~15 (\chref{c20}).
\end{warningbox}

\begin{insightbox}[Key takeaways -- three things to remember]
\begin{enumerate}[nosep, leftmargin=1.6em]
  \item \emph{A sum of distances reroots by a shift, not by a shortcut.} $\max$-type eccentricity
        obeys \thmref{c09-eccend}'s two-endpoint identity; $\sum$-type totals obey
        $f[v]=f[u]+S-2W_v$ -- one mass DFS plus one replay (\probref{c09-summit}).
  \item \emph{The size table \emph{is} the mass table.} The only new object versus CF~219D is
        summing $a$ instead of counting nodes -- the whole chapter's $\size$ discipline returns
        the value you need in one pass (\thmref{c09-onepass}, \secref{c09-radius}).
  \item \emph{Choose $\max$ vs $\sum$ before coding.} Best-worst-case root $\to$ centre / radius;
        best-total root $\to$ reroot shift. Naming the objective names the algorithm
        (\thmref{c09-centerrule} vs \probref{c09-summit}).
\end{enumerate}
\end{insightbox}

\begin{hardeningbox}
\begin{itemize}[nosep, leftmargin=1.4em]
  \item \textbf{Generalises}: any ``answer for every possible root'' question -- all-roots DP,
        tree hash re-rooting, sum of pairwise distances -- is one seed DFS plus one edge-transfer
        replay; only the shift formula changes per objective (Chapter~15's catalogue).
  \item \textbf{The mass was the whole difficulty}: a plain DFS already gives $\mathit{sz}$ and
        $f[\text{root}]$; promoting to all roots needs only the $\bigO{1}$ shift, which is why
        this is $\lv{5}$ for the composition, not the parts (\probref{c09-summit}).
  \item \textbf{Alternatives}: a centroid-decomposition pass computes distance \emph{counts} in a
        different $\bigO{n\log n}$ way (Chapter~24) but is overkill here; an explicit all-pairs
        BFS is the honest $\Theta(n^2)$ fallback for tiny $n$ and fails at $2\cdot10^{5}$
        (\coderef{c09-kit}).
\end{itemize}
\end{hardeningbox}

\chaptersection{Tool Kit Inventory}

The machine assembled so far -- cumulative through Chapter~12. Chapters~1--11's rows are
carried forward; Chapter~12's rows are appended. Pointers only, never prose.

\begin{center}
\small
\renewcommand{\arraystretch}{1.25}
\begin{tabular}{@{}p{0.24\linewidth} l c p{0.27\linewidth} p{0.16\linewidth}@{}}
\toprule
\textbf{Tool} & \textbf{Since} & \textbf{Level} & \textbf{Trigger phrase} & \textbf{Where} \\
\midrule
\multicolumn{5}{l}{\itshape Chapters 1--11 -- carried forward} \\
Binary-tree height/size recursion & 1 & \lv1 & ``compute at the node from its two children'' & \chref{c01} \\
Bottom-up DP pattern & 2 & \lv2 & ``a node folds its children's answers'' & \chref{c02} \\
Parent-guarded rooting / representation & 4 & \lv2 & ``root a general tree with a parent parameter'' & \chref{c04} \\
DFS/BFS visit order & 3 & \lv2 & ``preorder, postorder, level order'' & \coderef{c03-dfs} \\
\midrule
\multicolumn{5}{l}{\itshape Chapter 12 -- new this chapter} \\
Measures DFS ($\depth/\height/\size$) & 9.1 & \lv2 & ``how deep / how tall / how many below'' & \coderef{c09-dfs3} \\
Iterative two-phase DFS & 9.1 & \lv3 & ``the chain may be $5\cdot10^{5}$ deep'' & \coderef{c09-iter} \\
Two-sweep diameter & 9.2 & \lv3 & ``longest path, unweighted, in one shot'' & \coderef{c09-sweep} \\
DP diameter (join two arms) & 9.2 & \lv3 & ``diameter that must carry weights / a predicate'' & \thmref{c09-dpjoin}, \coderef{c09-dp} \\
Centroid via maxpart & 9.3 & \lv3 & ``a safe node to divide-and-conquer on'' & \thmref{c09-centfind}, \coderef{c09-cent} \\
Eccentricity from two ends & 9.4 & \lv4 & ``worst-case distance to the farthest node'' & \thmref{c09-eccend} \\
Radius / centre rule & 9.4 & \lv3 & ``best root by maximum distance; $R=\lceil D/2\rceil$'' & \thmref{c09-centerrule} \\
Distance-sum reroot shift & 9.4 & \lv4 & ``a $\sum$ over all roots, one replay'' & \coderef{c09-eccsweep} \\
Weighted all-roots reroot (Summit) & Summit & \lv5 & ``$\max$/$\min$ of $\sum\mathit{dist}\cdot a$ for every root'' & \coderef{c09-summit} \\
\bottomrule
\end{tabular}
\end{center}

\chaptersection{Related Topics}

\begin{itemize}[nosep, leftmargin=1.4em]
  \item \chref{c01}, \chref{c02} -- backward: this chapter's ``depth down, height/size up'' is
        Ch.~1's binary recursion (\thmref{c01-height-recursion}) re-framed over a general child
        list with a parent guard; the folding pattern is Ch.~2's, cited not re-taught
        (\defref{c09-measures}, \thmref{c09-onepass}).
  \item \chref{c04} -- backward: rooting by an explicit parent parameter, and the LC~310
        leaf-peelling centre, are Chapter~4's (\secref{c09-radius} keeps only the distance side).
  \item \chref{c10} -- forward: the ``take it or leave it per child'' tree-DP template that
        CSES~1130 previews is Chapter~13's subject; this chapter only hands it a size table.
  \item \chref{c11} -- forward: per-pair \emph{query} distance $d(u,v)$ after preprocessing is
        LCA, Chapter~16; here every metric is one-shot. \coderef{c09-iter} is the stack-safe
        DFS that chapter's kit reuses.
  \item \chref{c18} -- forward: centroid \emph{decomposition} -- recursing the $n/2$ split to
        process all paths -- is Chapter~24's; \thmref{c09-centexist}'s walk-to-heavy proof is
        its blueprint, and CF~342E is gated behind it (\secref{c09-centroid}).
  \item \chref{c20} -- forward: rerooting \emph{as a theory} (any edge-transfer rule, not just
        the two here) is Chapter~15's; \probref{c09-capital} and \probref{c09-summit} are its
        two worked special cases.
\end{itemize}

\begin{lastword}[Last Word -- Chapter 12]
\emph{``Chapters~1--11 built trees to \textbf{store} things -- values, keys, order. This chapter
measured a tree that already exists: how deep it reaches, how far apart its ends are, and where
it balances.''}

\medskip
The whole chapter is one discipline repeated five times: \emph{scan the tree a constant number of
times, then read a number off each node}. A parent-guarded DFS returns depth, height and size
together (\thmref{c09-onepass}); call it twice and the farthest-node lemma (\lemref{c09-sweepend})
hands you the diameter (\thmref{c09-twosweep}); re-read the same size table and $\mathit{maxpart}$
falls out with the single upward piece $n-\size(u)$ separating a solved centroid from a wrong
answer (\thmref{c09-centfind}); keep the two sweeps' distance arrays and every eccentricity is a
$\max$ over the diameter ends (\thmref{c09-eccend}); and replace that $\max$ with a $\sum$ and a
signed edge-transfer gives all roots' totals in one replay (\probref{c09-summit}). The centroid is
a middle by \emph{count}, the centre a middle by \emph{distance} -- $T$ conflates them into the pair
$\{1,3\}$, the broom never does.

\medskip
Carry three reflexes forward:
\begin{enumerate}[nosep, leftmargin=1.6em]
  \item \emph{Never pay per pair for a one-shot metric.} Two BFS sweeps, not $\binom{n}{2}$
        distances (\thmref{c09-twosweep}); one reroot replay, not $n$ DFS passes (\probref{c09-summit}).
  \item \emph{The forgotten term is always the parent's side.} $\mathit{maxpart}$ without
        $n-\size(u)$, an eccentricity from one end, a size check that misses $\mathit{sz}[r]=n$ --
        the upward direction is where this chapter's bugs hide (\defref{c09-centroid},
        \thmref{c09-eccend}).
  \item \emph{Name the objective before choosing the tool.} ``Farthest'' $\to$ radius/centre;
        ``how many'' $\to$ centroid; ``sum over all'' $\to$ reroot -- the deciding word is in the
        problem text.
\end{enumerate}

\medskip
And three habits, drilled until boring: write the measures DFS once and let a parameter choose
the root; reach for \coderef{c09-iter} the moment $n$ clears $10^{5}$; and when a metric asks for
\emph{every} root, hunt for the two numbers that describe crossing one edge before you write a
second traversal (\coderef{c09-kit}).

\medskip
\textbf{Next:} \chref{c10} takes the size table you now trust and turns it into an \emph{optimisation
engine}: tree dynamic programming, where each node computes an answer from its children's answers
over a fixed recurrence. The folding discipline of \secref{c09-measures} becomes a design pattern,
the centroid's $n/2$ split becomes a divide-and-conquer budget, and the reroot shift you met at the
Summit graduates into a full technique (Chapter~15). The metrics are done -- from here, trees start
\emph{computing}.
\end{lastword}
